\section{Exceptions}
\begin{verbatim}
try { // do something }
catch (Exception e) { // handle exception }
finally { // clean up
   // no matter exception occurred or not 
}
\end{verbatim}
\begin{center}
  \vspace{-5mm}
  \includegraphics[width=0.5\linewidth]{exception}
  \vspace{-4mm}
\end{center}
\textcolor{Blue}{Checked exceptions} might still happen with perfect code,
programmer has no control.  Must be handled to compile.\\
\textcolor{Blue}{Unchecked exceptions} are caused by programmer errors that cause
Runtime Exceptions.
\paragraph{Catch}
\begin{verbatim}
} catch (ExceptionX e) { 
   ...
} catch (ExceptionY e) {
\end{verbatim}
Exceptions are checked in order. If $\verb|ExceptionY|<:\verb|ExceptionX|$, the
second catch will never run $\Rightarrow$ \textcolor{Bittersweet}{Exception!}
\paragraph{Design principles}
\begin{enumerate}
  \item Catch exceptions to clean up, even if just ``passing the buck''
  \item Do NOT catch all (Exception e)
  \item Do not overreact (Exiting the program)
  \item Do not break abstraction (Leaking abstraction throwing \verb|FileNotFound| etc.)
  \item Do not use as flow control mechanism
\end{enumerate}
